\section{Тема 11. Аналитическая геометрия}

\subsection{Постановка задачи (вариант~4)}

Заданы уравнения двух пересекающихся прямых на плоскости:
$a_1x+b_1y+c_1=0$ и $a_2x+b_2y+c_2=0$.
Найти (в градусах и минутах) угол между ними.

\subsection{Математическое обоснование}

Вектор нормали к прямой $a x+by+c=0$ есть $\vec n=(a,b)$.
Угол между прямыми равен углу между их нормалями (или его
дополнению до $180°$, поэтому берём острый угол):
\[
  \cos\varphi = \frac{|\vec n_1 \cdot \vec n_2|}{|\vec n_1|\,|\vec n_2|}
  = \frac{|a_1a_2+b_1b_2|}{\sqrt{a_1^2+b_1^2}\sqrt{a_2^2+b_2^2}}.
\]
Перевод радиан в градусы и минуты:
\[
  \varphi^{\circ} = \lfloor \varphi \cdot 180/\pi \rfloor,\qquad
  \varphi' = \operatorname{round}\bigl((\varphi\cdot180/\pi - \varphi^\circ)\cdot60\bigr).
\]

\subsection{Блок-схема алгоритма}

\begin{center}
\begin{tikzpicture}[
  node distance=10mm, start chain=going below,
  nodes={draw=black, top color=white, bottom color=lightgray!50,
         minimum width=5.6cm, align=center, font=\small},
  arrow/.style={->, >=Stealth, thick},
]
  \node[draw, rounded corners, on chain] {Start};
  \node[shape=trapezium, trapezium left angle=70,
        trapezium right angle=110, on chain, draw, join=by arrow]
       {$a_1,b_1,c_1,\ a_2,b_2,c_2$};
  \node[shape=rectangle, on chain, draw, join=by arrow]
       {$d := a_1a_2+b_1b_2$};
  \node[shape=rectangle, on chain, draw, join=by arrow]
       {$n_1 := \sqrt{a_1^2{+}b_1^2}$,\quad $n_2 := \sqrt{a_2^2{+}b_2^2}$};
  \node[shape=rectangle, on chain, draw, join=by arrow]
       {$\cos\varphi := |d|/(n_1 n_2)$};
  \node[shape=rectangle, on chain, draw, join=by arrow]
       {$\varphi := \arccos(\cos\varphi)$};
  \node[shape=rectangle, on chain, draw, join=by arrow]
       {${}^\circ,\ '$ из $\varphi$};
  \node[shape=trapezium, trapezium left angle=70,
        trapezium right angle=110, on chain, draw, join=by arrow]
       {$\varphi$};
  \node[draw, rounded corners, on chain, join=by arrow] {End};
\end{tikzpicture}
\end{center}

\subsection{Текст программы}

\lstinputlisting[style=Cstyle]{../src/t11.c}

\subsection{Тестовые данные}

\begin{center}
\begin{tabular}{l|l}
\toprule
Прямые & Угол \\
\midrule
$3x+4y=0$, $4x-3y=0$ & $90°\,0'$ (перпендикулярны) \\
$x+y=0$, $x-y=0$ & $90°\,0'$ \\
\bottomrule
\end{tabular}
\end{center}

\textbf{Результат выполнения программы:}
\begin{lstlisting}[language={},basicstyle=\ttfamily\small,frame=single]
Прямая 1 (a1 b1 c1): 3 4 0
Прямая 2 (a2 b2 c2): 4 -3 0
Угол между прямыми: 1.570796 рад = 90° 0'
\end{lstlisting}
